\begin{table}[h]
\centering
\caption{Effect of Exposure to Black Residents and Residential Zoning on Highway Construction}
\label{tab:results/suitabilityxblackshare+pctres}
\begin{threeparttable}
\begin{tabular*}{1.0\textwidth}{@{\extracolsep{\fill}}l*{1}{r}}
\toprule
 & Coefficient \\
\midrule
Exposure to Black Population & \makecell[tr]{-0.283 \\ (0.368)} \\
Est. Highway Suitability & \makecell[tr]{1.553{***} \\ (0.167)} \\
Exposure to Black Population x Est. Highway Suitability & \makecell[tr]{0.315{*} \\ (0.132)} \\
Percent Residential & \makecell[tr]{-1.285{**} \\ (0.415)} \\
R-squared & 0.273 \\
Observations & 5236 \\
\bottomrule
\end{tabular*}
\begin{tablenotes}[flushleft]
\footnotesize
\item Estimates from a Firth's bias-reduced logistic regression model of an indicator for highway construction from 1940-1959 on the covariates listed, as well as controls for housing, geographic, and demographic characteristics. The sample is restricted to squares that are outside the highway corridor. Percent Residential is the proportion of a square's area that is zoned for residential use. Highway Suitability is a logit value estimated for each square by a convolutional neural network based on geographical and geological factors. Exposure to Black Population is a location-specific weighted average of the share of Black residents in each square. Weights are computed by an exponential decay function of the straight-line distance between squares and set to zero beyond a distance of 3,000 meters. This measure is normalized within each city to account for city-level differences in Black population and concentration. Standard errors, reported in parenthesis, are \textcite{conley_gmm_1999} spatial HAC standard errors with a 1,500-meter distance cutoff.  *** p < 0.01, ** p < 0.05, * p < 0.1
\end{tablenotes}
\end{threeparttable}
\end{table}
